% GENERATED FILE. Edit canonical lessons, then run npm run generate:books.
% canonical-source-hash: fnv1a64:858ce063fca5509f
% canonical-lessons: MR-C387-vagalne, MR-C387-samavista-karne, MR-C387-visambne, MR-C387-tharne, MR-C387-patne

\chapter{What Counts}
\label{ch:mr-what-counts}
\hlchaptermodality{387}

\begin{chapteropening}
\textbf{By the end of this chapter:} \emph{I can say what is included, what is left out and what convinced me.}

Complete the last lesson of chapter 387: I can say what is included, what is left out and what convinced me.
\end{chapteropening}

\section[\texorpdfstring{\mr{वगळणे} (vagaḷṇe)}{वगळणे (vagaḷṇe)}]{\mr{वगळणे} (vagaḷṇe) — to leave out; to omit}

\label{lesson:MR-C387-vagalne}



\begin{quote}
Before the new one: say the Marathi for to take notice, then the Marathi for a loophole.
\end{quote}

\subsection*{You'll want to know: \mr{वगळणे}}
\textbf{\mr{वगळणे}} — \emph{vagaḷṇe} — \textquotedblleft{}to leave out; to omit\textquotedblright{}.

\begin{quote}
\mr{अहवालातून} \mr{त्याने} \mr{महत्त्वाचा} \mr{भाग} \mr{वगळला}.
\end{quote}

\begin{quote}
\emph{ahvālātūn tyāne mahattvācā bhāg vagaḷlā.}
\end{quote}

\begin{quote}
He left an important part out of the report.
\end{quote}

Past, with -\mr{ने} (he …): \mr{त्याने} \mr{वगळलं}, \emph{tyāne vagaḷlaṁ}.

The source takes -\mr{तून}; \mr{वगळता}, except.

\begin{cousinweb}
\textbf{\mr{वगळणे}} (\emph{vagaḷṇe}) is a tadbhava: CDIAL lists Marathi \emph{vagaḷṇẽ}, \textquotedblleft{}to drop disregardfully, cast out, omit\textquotedblright{}, under Sanskrit \emph{avagalati}, \textquotedblleft{}falls down, slips down\textquotedblright{} (CDIAL 742).
\end{cousinweb}

\begin{grammarlens}[title={What you've built}]
One more word for saying what happens if.
\end{grammarlens}

\subsection*{Your turn}
\begin{itemize}
\raggedright
  \item \emph{Say it:} \emph{vagaḷṇe}
  \item \emph{Say it:} \emph{vagaḷṇe}, once more
  \item \emph{Say it:} \emph{ahvālātūn tyāne mahattvācā bhāg vagaḷlā.} — He left an important part out of the report.
  \item \emph{Recall:} say the Marathi for to take notice, then the Marathi for a loophole, then say \emph{vagaḷṇe} again
\end{itemize}

\begin{tcolorbox}[breakable,colback=teal!4,colframe=teal!35!black,title={Before you move on}]
What does \mr{वगळणे} mean? (\textquotedblleft{}To leave out; to omit\textquotedblright{}.) Say it once more.
\end{tcolorbox}
\section[\texorpdfstring{\mr{समाविष्ट} \mr{करणे} (samāviṣṭa karṇe)}{समाविष्ट करणे (samāviṣṭa karṇe)}]{\mr{समाविष्ट} \mr{करणे} (samāviṣṭa karṇe) — to include}

\label{lesson:MR-C387-samavista-karne}



\begin{quote}
Before the new one: say the Marathi for a loophole, then the Marathi for to leave out.
\end{quote}

\subsection*{You'll want to know: \mr{समाविष्ट} \mr{करणे}}
\textbf{\mr{समाविष्ट} \mr{करणे}} — \emph{samāviṣṭa karṇe} — \textquotedblleft{}to include\textquotedblright{}.

\begin{quote}
\mr{नव्या} \mr{अभ्यासक्रमात} \mr{पर्यावरणाचा} \mr{विषय} \mr{समाविष्ट} \mr{केला} \mr{आहे}.
\end{quote}

\begin{quote}
\emph{navyā abhyāsakramāt paryāvarṇācā viṣay samāviṣṭa kelā āhe.}
\end{quote}

\begin{quote}
The environment has been included in the new syllabus.
\end{quote}

Past, with -\mr{ने} (he …): \mr{त्याने} \mr{समाविष्ट} \mr{केलं}, \emph{tyāne samāviṣṭa kelaṁ}.

Where it goes takes -\mr{त}: \mr{अभ्यासक्रमात} \mr{समाविष्ट}.

\begin{cousinweb}
\textbf{\mr{समाविष्ट} \mr{करणे}} (\emph{samāviṣṭa karṇe}) joins the tatsama \mr{समाविष्ट}, Sanskrit \emph{samāviṣṭa}, \textquotedblleft{}entered, filled\textquotedblright{} (Monier-Williams), to \mr{करणे}, \textquotedblleft{}to make\textquotedblright{}.
\end{cousinweb}

\begin{grammarlens}[title={What you've built}]
One more word for saying what happens if.
\end{grammarlens}

\subsection*{Your turn}
\begin{itemize}
\raggedright
  \item \emph{Say it:} \emph{samāviṣṭa karṇe}
  \item \emph{Say it:} \emph{samāviṣṭa karṇe}, once more
  \item \emph{Say it:} \emph{navyā abhyāsakramāt paryāvarṇācā viṣay samāviṣṭa kelā āhe.} — The environment has been included in the new syllabus.
  \item \emph{Recall:} say the Marathi for a loophole, then the Marathi for to leave out, then say \emph{samāviṣṭa karṇe} again
\end{itemize}

\begin{tcolorbox}[breakable,colback=teal!4,colframe=teal!35!black,title={Before you move on}]
What does \mr{समाविष्ट} \mr{करणे} mean? (\textquotedblleft{}To include\textquotedblright{}.) Say it once more.
\end{tcolorbox}
\section[\texorpdfstring{\mr{विसंबणे} (visambṇe)}{विसंबणे (visambṇe)}]{\mr{विसंबणे} (visambṇe) — to rely on}

\label{lesson:MR-C387-visambne}



\begin{quote}
Before the new one: say the Marathi for to leave out, then the Marathi for to include.
\end{quote}

\subsection*{You'll want to know: \mr{विसंबणे}}
\textbf{\mr{विसंबणे}} — \emph{visambṇe} — \textquotedblleft{}to rely on\textquotedblright{}.

\begin{quote}
\mr{मी} \mr{फक्त} \mr{स्वतःवर} \mr{विसंबतो}.
\end{quote}

\begin{quote}
\emph{mī phakta svataḥvar visambto.}
\end{quote}

\begin{quote}
I rely only on myself.
\end{quote}

Past (he …): \mr{तो} \mr{विसंबला}, \emph{to visamblā}.

The one relied on takes -\mr{वर}; \mr{विसंबून} \mr{राहणे}, to depend.

\begin{cousinweb}
\textbf{\mr{विसंबणे}} (\emph{visambṇe}) is inherited: CDIAL lists Marathi \emph{visãbṇẽ}, \textquotedblleft{}to pause, take rest\textquotedblright{}, under Sanskrit \emph{viśrāmyati}, \textquotedblleft{}takes rest\textquotedblright{} (CDIAL 11959); today it means to rest one's weight on someone.
\end{cousinweb}

\begin{grammarlens}[title={What you've built}]
One more word for saying what happens if.
\end{grammarlens}

\subsection*{Your turn}
\begin{itemize}
\raggedright
  \item \emph{Say it:} \emph{visambṇe}
  \item \emph{Say it:} \emph{visambṇe}, once more
  \item \emph{Say it:} \emph{mī phakta svataḥvar visambto.} — I rely only on myself.
  \item \emph{Recall:} say the Marathi for to leave out, then the Marathi for to include, then say \emph{visambṇe} again
\end{itemize}

\begin{tcolorbox}[breakable,colback=teal!4,colframe=teal!35!black,title={Before you move on}]
What does \mr{विसंबणे} mean? (\textquotedblleft{}To rely on\textquotedblright{}.) Say it once more.
\end{tcolorbox}
\section[\texorpdfstring{\mr{ठरणे} (ṭharṇe)}{ठरणे (ṭharṇe)}]{\mr{ठरणे} (ṭharṇe) — to turn out; to be settled}

\label{lesson:MR-C387-tharne}



\begin{quote}
Before the new one: say the Marathi for to include, then the Marathi for to rely on.
\end{quote}

\subsection*{You'll want to know: \mr{ठरणे}}
\textbf{\mr{ठरणे}} — \emph{ṭharṇe} — \textquotedblleft{}to turn out; to be settled\textquotedblright{}.

\begin{quote}
\mr{त्याचा} \mr{निर्णय} \mr{योग्य} \mr{ठरला}.
\end{quote}

\begin{quote}
\emph{tyācā nirṇay yogya ṭharlā.}
\end{quote}

\begin{quote}
His decision turned out right.
\end{quote}

Past (he …): \mr{तो} \mr{ठरला}, \emph{to ṭharlā}.

Also to be fixed: \mr{लग्न} \mr{ठरलं}; \mr{ठरवणे}, to decide, is its causative.

\begin{cousinweb}
\textbf{\mr{ठरणे}} (\emph{ṭharṇe}) is inherited: CDIAL lists Marathi \emph{ṭharṇẽ}, \textquotedblleft{}to be settled\textquotedblright{}, under \emph{*stharati}, \textquotedblleft{}is firm\textquotedblright{} (CDIAL 13742).
\end{cousinweb}

\begin{grammarlens}[title={What you've built}]
One more word for saying what happens if.
\end{grammarlens}

\subsection*{Your turn}
\begin{itemize}
\raggedright
  \item \emph{Say it:} \emph{ṭharṇe}
  \item \emph{Say it:} \emph{ṭharṇe}, once more
  \item \emph{Say it:} \emph{tyācā nirṇay yogya ṭharlā.} — His decision turned out right.
  \item \emph{Recall:} say the Marathi for to include, then the Marathi for to rely on, then say \emph{ṭharṇe} again
\end{itemize}

\begin{tcolorbox}[breakable,colback=teal!4,colframe=teal!35!black,title={Before you move on}]
What does \mr{ठरणे} mean? (\textquotedblleft{}To turn out; to be settled\textquotedblright{}.) Say it once more.
\end{tcolorbox}
\section[\texorpdfstring{\mr{पटणे} (paṭṇe)}{पटणे (paṭṇe)}]{\mr{पटणे} (paṭṇe) — to be convinced; to make sense (to)}

\label{lesson:MR-C387-patne}



\begin{quote}
Before the new one: say the Marathi for to rely on, then the Marathi for to turn out.
\end{quote}

\subsection*{You'll want to know: \mr{पटणे}}
\textbf{\mr{पटणे}} — \emph{paṭṇe} — \textquotedblleft{}to be convinced; to make sense (to)\textquotedblright{}.

\begin{quote}
\mr{तुमचं} \mr{म्हणणं} \mr{मला} \mr{पटलं}.
\end{quote}

\begin{quote}
\emph{tumcaṁ mhaṇṇaṁ malā paṭlaṁ.}
\end{quote}

\begin{quote}
What you said convinced me.
\end{quote}

Past (to him …): \mr{त्याला} \mr{पटलं}, \emph{tyālā paṭlaṁ}.

Dative: \mr{मला} \mr{पटलं}, I am convinced; \mr{पटवणे}, to convince, is its causative.

\begin{cousinweb}
\textbf{\mr{पटणे}} (\emph{paṭṇe}) is a Marathi verb whose origin is not settled; \mr{पटवणे}, to convince, is built on it.
\end{cousinweb}

\begin{grammarlens}[title={What you've built}]
One more word for saying what happens if.
\end{grammarlens}

\subsection*{Your turn}
\begin{itemize}
\raggedright
  \item \emph{Say it:} \emph{paṭṇe}
  \item \emph{Say it:} \emph{paṭṇe}, once more
  \item \emph{Say it:} \emph{tumcaṁ mhaṇṇaṁ malā paṭlaṁ.} — What you said convinced me.
  \item \emph{Recall:} say the Marathi for to rely on, then the Marathi for to turn out, then say \emph{paṭṇe} again
\end{itemize}

\begin{tcolorbox}[breakable,colback=teal!4,colframe=teal!35!black,title={Before you move on}]
What does \mr{पटणे} mean? (\textquotedblleft{}To be convinced; to make sense (to)\textquotedblright{}.) Say it once more.
\end{tcolorbox}
